% !TEX root = task-2-12-3.tex
% Editable task bank. Append new BankTask blocks at the end of the document.
\documentclass[12pt,a4paper]{article}
\usepackage{physics-textbook}
\usepackage{graphicx,needspace}
\usetikzlibrary{decorations.pathmorphing}
\newif\ifShowSolutions
\ifdefined\TutorEdition
  \ShowSolutionstrue
\else
  \ShowSolutionsfalse % Change to \ShowSolutionstrue to display solutions.
\fi
\newcounter{banktask}
\newsavebox{\BankTaskBox}
\newcommand{\BankTaskContent}[3]{%
  \noindent\textbf{\thebanktask.}\quad #2\par
  \ifShowSolutions
    \Needspace{50mm}% Keep the solution heading with its opening text or diagram.
    #3\par
    \ifdefined\TutorEdition
      \ifcsname TutorSource#1\endcsname
        \nopagebreak\smallskip{\small\csname TutorSource#1\endcsname\par}%
      \fi
    \fi
  \fi
}
\newcommand{\BankTask}[3]{%
  \refstepcounter{banktask}\par\medskip
  \sbox{\BankTaskBox}{\begin{minipage}{\linewidth}%
    \BankTaskContent{#1}{#2}{#3}\end{minipage}}%
  \ifdim\dimexpr\ht\BankTaskBox+\dp\BankTaskBox\relax>\dimexpr\textheight-20mm\relax
    \Needspace{35mm}\BankTaskContent{#1}{#2}{#3}%
  \else
    \Needspace{\dimexpr\ht\BankTaskBox+\dp\BankTaskBox\relax}%
    \noindent\usebox{\BankTaskBox}\par
  \fi
  \medskip
}
\begin{document}
\BookSubsection{2.12.3}{Тепловая машина Карно}


% Source sdamgia-1121: https://phys-ege.sdamgia.ru/problem?id=1121
\BankTask{sdamgia-1121}{
Температура нагревателя тепловой машины $900\,\text{К}$, температура холодильника на $300\,\text{К}$ меньше, чем у нагревателя. Каков максимально возможный КПД машины? ((В процентах, округлив до целых.))
}{
\smallskip\noindent\textbf{Решение}\par
Температура холодильника равна $T_{\text{x}}=T_{\text{н}} - 300\text{К}=900\text{К} - 300\text{К}=600\text{К}.$ Максимально возможный КПД тепловой машины равен КПД машины Карно: 

$\eta = \frac{T_{\text{н}} - T_{\text{x}}}{T_{\text{н}}} = \frac{900\text{К} - 600\text{К}}{900\text{К}} = \frac{1}{3} \approx33\%.$

Ответ: $33$.
}

% Source sdamgia-6818: https://phys-ege.sdamgia.ru/problem?id=6818
\BankTask{sdamgia-6818}{
В идеальной тепловой машине абсолютная температура нагревателя отличается от температуры холодильника в $2$ раза. Чему равен КПД этой машины? (В процентах.)
}{
\smallskip\noindent\textbf{Решение}\par
КПД идеальной тепловой машины вычисляется по формуле: $\eta= ( 1 - \frac{T_{\text{хол}}}{T_{\text{нагр}}} ) \cdot 100\%= ( 1 - \frac{1}{2} ) \cdot 100\%=50\%.$

Ответ: $50$.
}

% Source sdamgia-27978: https://phys-ege.sdamgia.ru/problem?id=27978
\BankTask{sdamgia-27978}{
Тепловой двигатель, работающий по циклу Карно, имеет нагреватель с температурой +$127^\circ$C и холодильник с температурой +$27^\circ$C. Какое количество теплоты этот двигатель отдает холодильнику за один цикл работы, если за каждый цикл он получает от нагревателя количество теплоты, равное $2\,\text{кДж}$?

(В кДж.)
}{
\smallskip\noindent\textbf{Решение}\par
КПД тепловой машины найдем по формулам $\eta= 1 - \frac{T_2}{T_1}$ и $\eta= 1 - \frac{Q_2}{Q_1} .$ Найдем количество теплоты, отданное холодильнику за цикл:

$Q_2= \frac{=T_2Q_1}{T_1} = \frac{300\text{К} \cdot 2\text{кДж}}{400\text{К}} =1{,}5\text{кДж}.$ 

Ответ: $1{,}5$.
}

% Source sdamgia-3617: https://phys-ege.sdamgia.ru/problem?id=3617
\BankTask{sdamgia-3617}{
В таблице приведена зависимость КПД идеальной тепловой машины от температуры ее нагревателя при неизменной температуре холодильника. Чему равна температура холодильника этой тепловой машины? (В кельвинах.)

\begin{center}\begin{tabular}{|c|c|c|c|c|c|}\hline
$T_{\text{н}}$, К & $400$ & $500$ & $600$ & $800$ & $1000$\\\hline
$\eta$, \% & $25$ & $40$ & $50$ & $62{,}5$ & $70$\\\hline\end{tabular}\end{center}
}{
\smallskip\noindent\textbf{Решение}\par
КПД идеальной машины связан с температурами нагревателя и холодильника соотношением $\eta= ( 1 - \frac{T_{\text{Х}}}{T_{\text{Н}}} ) \cdot 100\%.$ Отсюда для температуры холодильника имеем: $T_{\text{Х}}=T_{\text{Н}} ( 1 - \frac{\eta}{100\%} ) .$ Используя любой столбик из таблицы, получаем значение температуры $T_{\text{Х}}=400\text{К} \cdot ( 1 - \frac{25\%}{100\%} ) =300\text{К}.$ 

В качестве проверки можно рассчитать температуру холожильника еще для одной точки: $T_{\text{Х}}=1000\text{К} \cdot ( 1 - \frac{70\%}{100\%} ) =300\text{К}.$ Результаты совпадают, это хорошо.

Ответ: $300$.
}

% Source sdamgia-3291: https://phys-ege.sdamgia.ru/problem?id=3291
\BankTask{sdamgia-3291}{
У теплового двигателя, работающего по циклу Карно, температура нагревателя $500$ К, а температура холодильника $300$ К. Рабочее тело за один цикл получает от нагревателя количество теплоты, равное $40\,\text{кДж}$. Какую работу совершает за цикл рабочее тело двигателя? Ответ укажите в килоджоулях.
}{
\smallskip\noindent\textbf{Решение}\par
КПД теплового двигателя равен $\eta= \frac{A}{Q_{\,\text{н}}} ,$ КПД идеальной тепловой машины $\eta= \frac{T_{\,\text{н}} - T_x}{T_{\,\text{н}}} .$ Приравняем  формулы: 

\[\frac{T_{\,\text{н}} - T_x}{T_{\text{н}}}= \frac{A}{Q} ,\]
 отсюда для работы, которую совершает рабочее тело за один цикл, получаем: 

\[A= \frac{Q ( T_{\,\text{н}} - T_x )}{T_{\,\text{н}}} = \frac{40\,\text{кДж} \cdot ( 500K - 300K )}{500K} \approx 16\,\text{кДж}.\]

Ответ: $16\,\text{кДж}$.
}

% Source sdamgia-1220: https://phys-ege.sdamgia.ru/problem?id=1220
\BankTask{sdamgia-1220}{
Температура нагревателя идеального теплового двигателя Карно равна $527 ^\circ\text{C} ,$ а температура холодильника равна $27 ^\circ\text{C} .$ Рабочее тело получает от нагревателя за один цикл количество теплоты $25\,\text{кДж}$. Какую работу совершает за цикл рабочее тело двигателя?
}{
\smallskip\noindent\textbf{Решение}\par
Найдем значения температур нагревателя и холодильника в абсолютной шкале температур: 

$T_{\text{н}}=273K + 527K=800K,$

$T_{\text{x}}=273K + 27K=300K.$

КПД теплового двигателя $\eta= \frac{A}{Q_{\text{н}}} ,$ КПД идеальной тепловой машины $\eta= \frac{T_{\text{н}} - T_x}{T_{\text{н}}}$

Приравняем две формулы для КПД цикла Карно: 

$\frac{T_{\text{н}} - T_{\text{x}}}{T_{\text{н}}} = \frac{A}{Q} ,$ 

отсюда для работы, которую совершает рабочее тело за один цикл, имеем: 

$A= \frac{Q ( T_{\text{н}} - T_{\text{x}} )}{T_{\text{н}}} = \frac{25\text{кДж} \cdot 500\text{К}}{800\text{К}} \approx 15{,}6\text{кДж}.$

Ответ: $15{,}6\,\text{кДж}$.
}

% Source sdamgia-7302: https://phys-ege.sdamgia.ru/problem?id=7302
\BankTask{sdamgia-7302}{
В тепловой машине, работающей по циклу Карно, газ совершает за один цикл работу $225\,\text{Дж}$. Температура нагревателя равна $327^\circ$C, температура холодильника равна $27^\circ$C. Определите количество теплоты, получаемое газом за один цикл.
}{
\smallskip\noindent\textbf{Решение}\par
КПД идеальной тепловой машины с нагревателем с температурой $T_1 = 273 + 327=600 K$ и холодильником с температурой $T_2= 273 + 27 = 300\text{К}$ равен:

$\eta = 1 - \frac{T_2}{T_1} = 1 - \frac{300}{600} = 0.5$

Так как КПД теплового двигателя $\eta= \frac{A}{Q_1} ,$ то количество теплоты, которое получает газ за один цикл, может быть вычислено по формуле:

$Q= \frac{A}{\eta} = \frac{225}{0{,}5} = 450 \text{Дж}.$

Ответ: $450\,\text{Дж}$.
}

% Source sdamgia-10484: https://phys-ege.sdamgia.ru/problem?id=10484
\BankTask{sdamgia-10484}{
Тепловая машина с максимально возможным КПД имеет в качестве нагревателя резервуар с водой, а в качестве холодильника — сосуд со льдом при $0^\circ$C. При совершении машиной работы $1\,\text{МДж}$ растаяло $12{,}1\,\text{кг}$ льда. Определите температуру воды в резервуаре.
}{
\smallskip\noindent\textbf{Решение}\par
При таянии лед получил от газа количество теплоты $Q_{\text{x}}=\lambda m=12{,}1\text{кг} \cdot 3{,}3 \cdot 10 ^{5} \text{Дж}/\text{кг} = 3{,}993\text{МДж}.$ 

КПД тепловой машины равен $\eta= \frac{A}{Q_{\text{н}}} ,$ причем количество теплоты, полученное газом от нагревателя $Q_{\text{н}}=A + Q_{\text{x}},$ откуда $\eta= \frac{A}{Q_{\text{x}} + A} .$ 

С другой стороны КПД идеальной тепловой машины $\eta=1 - \frac{T_{\text{x}}}{T_{\text{н}}} .$ Таким образом:

$T_{\text{н}}= \frac{T_{\text{x}}}{1 - \eta} = \frac{T_{\text{x}}}{1 - \frac{A}{Q_{\text{x}} + A}} =T_{\text{x}} ( 1 + \frac{A}{Q_{\text{x}}} ) =273 \cdot ( 1 + \frac{1}{3{,}993} ) \approx 341\text{К}.$

Ответ: $341\,\text{К}$.
}

% Source sdamgia-4359: https://phys-ege.sdamgia.ru/problem?id=4359
\BankTask{sdamgia-4359}{
Идеальная тепловая машина с температурой холодильника $300\,\text{К}$ и температурой нагревателя $400\,\text{К}$ за один цикл своей работы получает от нагревателя количество теплоты $10\,\text{Дж}$. За счет совершаемой машиной работы груз массой $10\,\text{кг}$ поднимается вверх с поверхности земли. На какую высоту над землей поднимется этот груз через $100$ циклов работы машины? (В метрах.)
}{
\smallskip\noindent\textbf{Решение}\par
Выпишем два выражения для КПД идеальной тепловой машины: $\eta = \frac{T_{\text{н}} - T_{\text{x}}}{T_{\text{н}}} ,$ $\eta = \frac{A}{Q_1} .$ Приравняв их, находим полезную работу за один цикл: $A=Q_1 \frac{T_{\text{н}} - T_{\text{x}}}{T_{\text{н}}} .$ Эта работа идет на изменение потенциальной энергии груза $\Delta E_{\text{п}}=mgh.$ 

Таким образом, за $100$ циклов груз поднимется на высоту: 

$H = 100 \cdot \frac{Q_1}{mg} \cdot \frac{T_{\text{н}} - T_{\text{x}}}{T_{\text{н}}} =100 \cdot \frac{10\text{Дж}}{10\text{кг} \cdot 10 \text{м}/\text{с} ^{2}} \cdot \frac{400\text{К} - 300 \text{К}}{400\text{К}} =2{,}5\text{м}$

Ответ: $2{,}5\,\text{м}$.
}

% Source sdamgia-10959: https://phys-ege.sdamgia.ru/problem?id=10959
\BankTask{sdamgia-10959}{
Тепловая машина за один цикл совершает работу $25\,\text{Дж}$ и отдает холодильнику количество теплоты $75\,\text{Дж}$. Температура нагревателя этой машины $600\,\text{К}$, а температура холодильника $300\,\text{К}$. Во сколько раз КПД идеальной тепловой машины, работающей при тех же температурах нагревателя и холодильника, больше КПД рассматриваемой тепловой машины?
}{
\smallskip\noindent\textbf{Решение}\par
Для КПД идеальной тепловой машины справедлива формула: 

$\eta= \frac{T_{\text{н}} - T_{\text{x}}}{T_{\text{н}}} =0{,}5.$

С другой стороны, КПД теплового двигателя равен $\eta= \frac{A}{Q_{\text{н}}} ,$ где количество теплоты, полученное газом от нагревателя $Q_{\text{н}}=A + Q_{\text{x}}.$ Таким образом, получаем:

$\eta= \frac{A}{A + Q_{\text{x}}} = \frac{25}{25 + 75} =0{,}25.$

Следовательно, КПД идеальной тепловой машины в $2$ раза больше КПД рассматриваемой тепловой машины.

Ответ: $2$.
}


% Source 0D600A: https://ege.fipi.ru/bank/questions.php?search=1&proj=BA1F39653304A5B041B656915DC36B38&qid=0D600A
\BankTask{0D600A}{
У идеального теплового двигателя Карно температура нагревателя равна \(227\,^{\circ}\mathrm{C}\), а температура холодильника равна \(27\,^{\circ}\mathrm{C}\). Определите КПД теплового двигателя (в $\%$).
}{
\smallskip\noindent\textbf{Решение от GPT}\par
Температуры резервуаров: \(T_1=227+273=500\,\text{К}\), \(T_2=27+273=300\,\text{К}\). КПД двигателя Карно:
\[
\eta=1-\frac{T_2}{T_1}=1-\frac{300}{500}=0{,}4=40\%.
\]
Ответ: \(40\%\).
}

% Source EC33AF: https://ege.fipi.ru/bank/questions.php?search=1&proj=BA1F39653304A5B041B656915DC36B38&qid=EC33AF
\BankTask{EC33AF}{
У идеального теплового двигателя Карно температура нагревателя \(400\,\text{К}\), температура холодильника \(300\,\text{К}\). За один цикл двигатель совершает работу \(16\,\text{кДж}\). Какое количество теплоты получает рабочее тело от нагревателя за один цикл (в $\text{кДж}$)?
}{
\smallskip\noindent\textbf{Решение от GPT}\par
КПД двигателя Карно:
\[
\eta=1-\frac{T_2}{T_1}=1-\frac{300}{400}=0{,}25.
\]
Из \(\eta=A/Q_1\) получаем:
\[
Q_1=\frac{A}{\eta}=\frac{16}{0{,}25}=64\,\text{кДж}.
\]
Ответ: \(64\,\text{кДж}\).
}

% Source D67CAF: https://ege.fipi.ru/bank/questions.php?search=1&proj=BA1F39653304A5B041B656915DC36B38&qid=D67CAF
\BankTask{D67CAF}{
Температура нагревателя идеального теплового двигателя, работающего по циклу Карно, равна $T_1$, а коэффициент полезного действия этого двигателя равен $\eta$. За цикл рабочее тело двигателя получает от нагревателя количество теплоты $Q_1$.

Установите соответствие между физическими величинами и формулами, по которым их можно рассчитать.

\TaskMatchingColumns[0.58\linewidth]{Физические величины}{
А) количество теплоты, отдаваемое рабочим телом двигателя за цикл\par
Б) температура холодильника
}{Формулы}{
1) $\dfrac{T_1}{1-\eta}$\par
2) $T_1(1-\eta)$\par
3) $Q_1(1-\eta)$\par
4) $Q_1\eta$
}

\begin{center}
\begin{tabular}{|c|c|}\hline
А & Б \\ \hline
\phantom{00} & \phantom{00} \\ \hline
\end{tabular}
\end{center}
}{
\smallskip\noindent\textbf{Решение от GPT}\par
Обозначим отданное холодильнику количество теплоты через $Q_2>0$. Из $\eta=\dfrac{Q_1-Q_2}{Q_1}$ следует $Q_2=Q_1(1-\eta)$: А — 3. Для цикла Карно $\eta=1-\dfrac{T_2}{T_1}$, поэтому $T_2=T_1(1-\eta)$: Б — 2.

Ответ: А — 3, Б — 2.
}

% Source 352D7e: https://ege.fipi.ru/bank/questions.php?search=1&proj=BA1F39653304A5B041B656915DC36B38&qid=352D7e
\BankTask{352D7e}{
В тепловом двигателе, работающем по циклу Карно, температуру нагревателя понизили. Температура холодильника и количество теплоты, отданное рабочим телом холодильнику за цикл, остались прежними. Как изменились при этом КПД теплового двигателя и работа теплового двигателя за цикл?

Для каждой величины определите соответствующий характер её изменения:
\begin{enumerate}
\item увеличилась
\item уменьшилась
\item не изменилась
\end{enumerate}

\TaskChangeTable{КПД теплового двигателя}{Работа теплового двигателя за цикл}
}{
\smallskip\noindent\textbf{Решение от GPT}\par
Обозначим абсолютные температуры нагревателя и холодильника через \(T_{\text{н}}\) и \(T_{\text{х}}\), а модуль теплоты, отданной холодильнику, через \(Q_{\text{х}}\). КПД цикла Карно
\[
\eta=1-\frac{T_{\text{х}}}{T_{\text{н}}}
\]
при понижении \(T_{\text{н}}\) уменьшается (\(2\)). Из \(Q_{\text{н}}/Q_{\text{х}}=T_{\text{н}}/T_{\text{х}}\) следует
\[
A=Q_{\text{н}}-Q_{\text{х}}=Q_{\text{х}}\left(\frac{T_{\text{н}}}{T_{\text{х}}}-1\right).
\]
При неизменных \(Q_{\text{х}}\) и \(T_{\text{х}}\) работа также уменьшается (\(2\)). Ответ: \(22\).
}

\end{document}
